\documentclass[a4paper,11pt]{article}
\usepackage[utf8]{inputenc}
\usepackage{geometry}
\usepackage{xcolor}
\usepackage{listings}
\usepackage{hyperref}
\usepackage{graphicx}
\usepackage{amsmath}
\usepackage{amssymb}
\usepackage{tcolorbox}
\usepackage{float}

% Geometria della pagina
\geometry{top=2.5cm, bottom=2.5cm, left=2.5cm, right=2.5cm}

% Colori custom
\definecolor{ibmblue}{RGB}{0, 45, 156}
\definecolor{ibmred}{RGB}{250, 75, 76}
\definecolor{codegreen}{rgb}{0,0.6,0}
\definecolor{codegray}{rgb}{0.5,0.5,0.5}
\definecolor{codepurple}{rgb}{0.58,0,0.82}
\definecolor{backcolour}{rgb}{0.95,0.95,0.92}

% Setup per il codice Python
\lstdefinestyle{mystyle}{
    backgroundcolor=\color{backcolour},   
    commentstyle=\color{codegreen},
    keywordstyle=\color{magenta},
    numberstyle=\tiny\color{codegray},
    stringstyle=\color{codepurple},
    basicstyle=\ttfamily\footnotesize,
    breakatwhitespace=false,         
    breaklines=true,                 
    captionpos=b,                    
    keepspaces=true,                 
    numbers=left,                    
    numbersep=5pt,                  
    showspaces=false,                
    showstringspaces=false,
    showtabs=false,                  
    tabsize=2
}
\lstset{style=mystyle}

% Box personalizzati
\newtcolorbox{studentbox}{colback=blue!5!white,colframe=blue!75!black,title=Student Activity}
\newtcolorbox{theorybox}{colback=gray!5!white,colframe=gray!50!black,title=Theory Recap}
\newtcolorbox{warningbox}{colback=yellow!10!white,colframe=orange!80!black,title=Important: The Speedup}
\newtcolorbox{composerbox}{colback=purple!5!white,colframe=purple!75!black,title=IBM Composer View (Q-Sphere)}
\newtcolorbox{mathbox}{colback=red!5!white,colframe=red!75!black,title=Step-by-Step: The Math of Amplitude Amplification}

\title{\textbf{Laboratory Guide 3.5: Grover's Algorithm}\\
\large Quantum Search and Amplitude Amplification}
\author{Course: Quantum Computing Foundation}
\date{}

\begin{document}

\maketitle

\section{Introduction and Objectives}
In Lab 3, we used an Oracle to discover a global property of a function (Constant vs. Balanced). In this laboratory, we tackle a much more practical problem: \textbf{Unstructured Search}. Imagine looking for a specific name in a phone book that is completely scrambled. We will build \textbf{Grover's Algorithm}, which uses quantum interference to find the needle in the haystack quadratically faster than any classical computer.

\vspace{0.3cm}
\noindent \textbf{Prerequisites:}
\begin{itemize}
    \item Lesson QC14: \textit{Quantum Algorithms (Grover & Shor)}
    \item Lab 3: \textit{Oracles and Phase Kickback}
\end{itemize}

\noindent \textbf{Learning Objectives:}
\begin{enumerate}
    \item Understand the two components of Grover's Algorithm: The Oracle and the Diffuser.
    \item Mathematically track "Inversion about the Mean" (Amplitude Amplification).
    \item Implement a 2-qubit search algorithm in Qiskit.
\end{enumerate}

\begin{warningbox}
\textbf{The Quadratic Speedup}\\
Classically, finding one specific item in an unsorted list of $N$ items requires evaluating the list $N/2$ times on average, and $N$ times in the worst case ($\mathcal{O}(N)$).
Grover's Algorithm finds the item in roughly $\sqrt{N}$ steps ($\mathcal{O}(\sqrt{N})$).
For a database of 1,000,000 items, a classical computer needs up to 1,000,000 checks. A quantum computer needs only about 1,000 checks!
\end{warningbox}

\newpage

\section{Part A: The Anatomy of a Quantum Search}

Grover's Algorithm always follows the same logical loop:
\begin{enumerate}
    \item \textbf{Initialize:} Create an equal superposition of all possible answers.
    \item \textbf{The Oracle:} Flip the phase (make it negative) \textit{only} for the correct answer.
    \item \textbf{The Diffuser:} Perform an "Inversion about the Mean" to shrink the wrong answers and amplify the correct one.
    \item \textbf{Repeat:} Steps 2 and 3 are repeated $\approx \sqrt{N}$ times.
\end{enumerate}

\begin{mathbox}
\textbf{The 2-Qubit Magic}\\
We will search a database of $N=4$ items (using 2 qubits). The winning item is $|11\rangle$. Let's trace the amplitudes.
\begin{enumerate}
    \item \textbf{Superposition (H gates):}
    $$|\psi_1\rangle = \frac{1}{2}|00\rangle + \frac{1}{2}|01\rangle + \frac{1}{2}|10\rangle + \frac{1}{2}|11\rangle$$
    \textit{The average (mean) amplitude is $\frac{1}{2}$. Probabilities are $25\%$ each.}
    
    \item \textbf{The Oracle (Marks $|11\rangle$):} 
    The Oracle applies a negative phase to the winner.
    $$|\psi_2\rangle = \frac{1}{2}|00\rangle + \frac{1}{2}|01\rangle + \frac{1}{2}|10\rangle - \frac{1}{2}|11\rangle$$
    \textit{The new mean amplitude drops to $\frac{1}{4}$ (because three are $+0.5$ and one is $-0.5$).}
    
    \item \textbf{The Diffuser (Inversion about the Mean):}
    The Diffuser reflects every amplitude around the new mean ($\frac{1}{4}$). The formula is: $\text{New} = 2 \times \text{Mean} - \text{Old}$.
    \begin{itemize}
        \item For wrong answers ($+0.5$): $2(\frac{1}{4}) - \frac{1}{2} = \frac{1}{2} - \frac{1}{2} = 0$.
        \item For the winner ($-0.5$): $2(\frac{1}{4}) - (-\frac{1}{2}) = \frac{1}{2} + \frac{1}{2} = 1$.
    \end{itemize}
    $$|\psi_3\rangle = 0|00\rangle + 0|01\rangle + 0|10\rangle + 1|11\rangle$$
\end{enumerate}
\textbf{Result:} In exactly 1 step, the probability of measuring the winner $|11\rangle$ went from 25\% to 100\%!
\end{mathbox}



\section{Part B: Building Grover (IBM Quantum Composer)}
\textit{Platform: IBM Quantum Composer}

\subsection{Exercise 1: Visualizing the Amplification}
We will build the Oracle to mark $|11\rangle$. A gate that applies a phase only when both qubits are 1 is the \textbf{Controlled-Z (CZ) gate}.

\begin{studentbox}
\textbf{Action:}
\begin{enumerate}
    \item Put $q_0$ and $q_1$ in superposition using \texttt{H} gates.
    \item \textbf{The Oracle:} Apply a \texttt{CZ} gate connecting $q_0$ and $q_1$. Look at the Q-Sphere: the state $|11\rangle$ is now Red (negative phase).
    \item \textbf{The Diffuser:} 
    \begin{itemize}
        \item Apply \texttt{H} gates to both qubits.
        \item Apply \texttt{X} gates to both qubits.
        \item Apply a \texttt{CZ} gate between them.
        \item Apply \texttt{X} gates to both qubits.
        \item Apply \texttt{H} gates to both qubits.
    \end{itemize}
    \item \textbf{Look at the Probabilities bar chart.}
\end{enumerate}
\textbf{Expected Effect:} The bar chart will show exactly 100\% probability on the state $|11\rangle$. The algorithm successfully amplified the marked state and suppressed the others.
\end{studentbox}

\newpage

\section{Part C: Coding Grover's Algorithm (Qiskit)}

\subsection{Exercise 2: Python Implementation}
Now we implement the 2-qubit Grover search in Qiskit to find the state $|11\rangle$.

\begin{tcolorbox}[colback=green!5!white,colframe=green!75!black,title=Coding Task: Grover for $|11\rangle$]
\begin{lstlisting}[language=Python]
from qiskit import QuantumCircuit
from qiskit_aer import AerSimulator
from qiskit.visualization import plot_histogram

# 1. Setup
qc = QuantumCircuit(2, 2)

# 2. Initialization: Equal Superposition
qc.h([0, 1])
qc.barrier()

# 3. THE ORACLE (Marks state |11>)
# A Controlled-Z gate flips the phase only if both qubits are 1
qc.cz(0, 1)
qc.barrier()

# 4. THE DIFFUSER (Amplitude Amplification)
# Step 4a: Apply H gates
qc.h([0, 1])
# Step 4b: Apply X gates (shifts the |00> state to |11>)
qc.x([0, 1])
# Step 4c: Apply CZ (marks the new |11> state)
qc.cz(0, 1)
# Step 4d: Undo X and H
qc.x([0, 1])
qc.h([0, 1])
qc.barrier()

# 5. Measurement
qc.measure([0, 1], [0, 1])

# 6. Execution
simulator = AerSimulator()
job = simulator.run(qc, shots=1000)
counts = job.result().get_counts()
print(f"Search Results: {counts}")
\end{lstlisting}
\textbf{Analysis:} The output will be 100\% \texttt{'11'}. You found the hidden data in just 1 query!
\end{tcolorbox}

\section{Part D: Assessment (Quiz)}

\textbf{Q1. The Speedup}\\
If you have an unsorted database of 10,000 items, roughly how many operations does Grover's Algorithm need to find the target?
\begin{itemize}
    \item[A)] 10,000
    \item[B)] 5,000
    \item[C)] 100
\end{itemize}

\vspace{0.3cm}

\textbf{Q2. The Division of Labor}\\
In Grover's Algorithm, what is the specific job of the \textbf{Diffuser}?
\begin{itemize}
    \item[A)] To "mark" the correct answer by flipping its phase to negative.
    \item[B)] To invert all amplitudes around the mean, which shrinks the wrong answers and boosts the marked correct answer.
    \item[C)] To put all qubits into an initial equal superposition.
\end{itemize}

\vspace{1cm}
\hrule
\vspace{0.2cm}
\footnotesize{\textbf{Solutions:} Q1: C ($\sqrt{10000} = 100$); Q2: B (The Oracle marks the answer, the Diffuser amplifies it).}

\end{document}